% S06-M7 (package abhijit/s06-math): Betti incarnation of c_q. Labels: prop:s06-betti, lem:s06-twobytwo, rem:s06-betti-kp.
\section{The Betti incarnation of \texorpdfstring{$\cc_q$}{c\_q}}

Throughout this section $l=p$, $\Gamma=\Gamma_0(M)$, $V=V_k\otimes\cO=\Sym^{k-2}(\cO^2)$, and $\T_k(M)$ (the fixed-weight algebra, not the $\Lam$-adic $\T(M)$ of the section on $\Lam$-adic decomposition) is the Hecke algebra of $\Gamma_0(M)$
acting on $H^1(\Gamma,V)$. From \pathtt{if:s03-amalgam} we use \emph{only} two facts: the operator $\Psi^*\Psi$ on $H^1(\Gamma,V)^2$ is
given by a $2\times2$ matrix
\[\Psi^*\Psi=\begin{pmatrix}q+1&-T'_q\\-T''_q&q+1\end{pmatrix},\qquad T'_q,T''_q\in\T_k(M)[1/q],\quad T'_qT''_q=q^{2-k}T_q^2,\]
so that $\det(\Psi^*\Psi)=(q+1)^2-q^{2-k}T_q^2$ (stage 03: $\{T'_q,T''_q\}=\{T_q,q^{2-k}T_q\}$, the assignment depending on the
normalisation of stage 03 D4 / \pathtt{stage-03.md} \S0.2.8, which we do not need), and the definition
$\mathrm{LR}_q(M,k)_\m=\coker(\Psi^*\Psi)_\m$. Since $q\in\cO^\times$, all entries lie in $\T_k(M)\otimes\cO$. Weights are even: $-1\in\Gamma$
acts on $V$ by $(-1)^{k-2}$, so $H^1(\Gamma,V)=0$ for odd $k$ ($p$ odd). Put $c:=q^{(k-2)/2}\in\Z$.

\begin{proposition}[Betti incarnation of $\cc_q$]\label{prop:s06-betti}
Let $l=p>k\ge4$ (the range of \pathtt{if:s03-amalgam}), $q\not\equiv-1\bmod p$, and let $\m$ be non-Eisenstein for $\T_k(M)$ acting on
$H^1(\Gamma_0(M),V_k)$, with $H^1(\Gamma_0(M),V_k)_\m$ free of rank two over $\T_k(M)_\m$. Then:
\begin{enumerate}[label=(\alph*)]
\item $\coker(\Psi^*\Psi)_\m\cong\big(\T_k(M)_\m/(\delta)\big)^{\oplus2}$, where $\delta:=\det(\Psi^*\Psi)=(q+1)^2-q^{2-k}T_q^2$.
  Since $p\nmid2(q+1)$, $(\delta)=(T_q-\epsilon(q+1)q^{(k-2)/2})$ for the unique sign $\epsilon$ with this element in $\m$.
\item $\T_k(N)_{\m}/(\fa+\fb)\cong\T_k(M)_\m/(T_q-\epsilon(q+1)q^{(k-2)/2})$ (Ribet's level-raising module, with Ihara's lemma).
\item If $\rk_\cO\T'_{k,\m}=1$, then $\ord_\varpi\cc_q(f)=\tfrac12\,\mathrm{length}_\cO\coker(\Psi^*\Psi)_\m$. In general only
  $\ord_\varpi\cc_q(f)\le\tfrac12\mathrm{length}$ holds.
\end{enumerate}
(a) \stproved{} (linear algebra, Lemma~\ref{lem:s06-twobytwo}). (b) \stpmc{} (Ihara's lemma in weight $k$ and the cohomological
congruence module: \cite[Lemma~3.2 and \S2]{Dia91}, \cite[\S3]{Dia89}). (c) \stpmc{} (from (a), (b)).
Here: in (a), if neither sign gives an element of $\m$, then $\delta\in\T_k(M)_\m^\times$ and both sides are $0$ for either sign;
in (b), $\T_k(N)_\m$ is the localisation of the weight-$k$ Hecke algebra of $\Gamma_0(N)$ (with $U_q$; $N=qM$) at the maximal ideal $\m_N$
generated by $\m$ (on the $T_\ell$, $\ell\nmid Np$, and $U_\ell$, $\ell\mid M$) and $U_q-\epsilon c$, and $\fa,\fb$ are its $q$-new and
$q$-old kernels; in (c), $f$ is a $q$-new eigenform in $\m_N$, $\T'_{k,\m}$ is the $q$-new quotient of $\T_k(N)_\m$, $\cc_q(f)$ is as in
Def.~\ref{def:s06-ctam}, and (H0) is assumed. \provisional{S06-M7: the degenerate case ``no sign'' in (a) and the maximal ideal
$\m_N$ of (b) are made explicit, and (H0) ($p\nmid\varphi(N)$, $q\not\equiv1$) is used in (b); statement otherwise unchanged.}
\end{proposition}

\begin{lemma}[$2\times2$ cokernels]\label{lem:s06-twobytwo}
Let $T$ be a commutative ring, $u\in T^\times$, $a,b\in T$, $X=\big(\begin{smallmatrix}u&-a\\-b&u\end{smallmatrix}\big)$ and $H$ a $T$-module. Then
$\coker(X:H^2\to H^2)\cong H/(u^2-ab)H$. \stproved{} \stformal{} for the case $H=T$ only (\lean{CommAlg.TwoByTwo.cokerEquivOfIsUnit}, \Tpath{CommAlg/TwoByTwo.lean}; the general $T$-module $H$ is not formalized).
\end{lemma}
\begin{proof}
With $E=\big(\begin{smallmatrix}1&0\\ bu^{-1}&1\end{smallmatrix}\big)$ and $E'=\big(\begin{smallmatrix}1&au^{-1}\\0&1\end{smallmatrix}\big)$ in $\GL_2(T)$,
$EXE'=\mathrm{diag}(u,\,u-abu^{-1})$. So $\coker X\cong\coker(u\cdot)\oplus\coker\big((u-abu^{-1})\cdot\big)=0\oplus H/(u^2-ab)H$.
\end{proof}

\begin{proof}[Proof of (a)]
Put $T:=\T_k(M)_\m$ and $H:=H^1(\Gamma,V)_\m\cong T^2$. Localisation at $\m$ is exact, so $\coker(\Psi^*\Psi)_\m$ is the cokernel of $\Psi^*\Psi$ on
$H^2$, which is the matrix above with entries in $T$. As $q\not\equiv-1\bmod p$, $u:=q+1\in\Zp^\times$. Lemma~\ref{lem:s06-twobytwo} gives
$\coker(\Psi^*\Psi)_\m\cong H/\delta H\cong(T/\delta)^2$, with $\delta=u^2-T'_qT''_q=(q+1)^2-q^{2-k}T_q^2$. Only $\det$ enters, so this is
independent of the $T'_q/T''_q$ choice of stage 03 D4. Next,
\[\delta=-q^{2-k}\big(T_q-(q+1)c\big)\big(T_q+(q+1)c\big),\]
and the two factors differ by $2(q+1)c\in\Zp^\times$ ($p\nmid2q(q+1)$). In the local ring $T$ at most one of them lies in $\m$; the other is a
unit. If $T_q-\epsilon(q+1)c\in\m$, then $(\delta)=(T_q-\epsilon(q+1)c)$; if neither lies in $\m$, $\delta$ is a unit.
\end{proof}

\noindent\emph{Cited inputs for (b).}
(B1) \emph{Ihara's lemma in weight $k$} \cite[Lemma~3.2, p.~210]{Dia91}
\lit{asterisque\_1991\_courbes-modulaires-et-courbes-de-shimura\_vol}{Diamond, ``Congruence primes for cusp forms of weight
$k\ge2$'', pp.~205--213: Lemma 3.2, p.~210 (PDF p.~214)} \stcited: for $p\nmid M(k-2)!$, the map
$\alpha:W(K/\cO)^2\to V_N(K/\cO)$ induced by the two degeneracy maps is injective, where $W$, $V_N$ are the images of integral parabolic
cohomology of $\Gamma_1(M)$, $\Gamma_1(M)\cap\Gamma_0(q)$ with coefficients $L_{k-2}$. Here $p>k$ and $p\nmid M$, so $p\nmid M(k-2)!$.
(B2) \emph{The cohomological congruence module} \cite[\S2, p.~207--208]{Dia91} (the weight-$k$ version of \cite[Prop.~3.3, p.~58--59]{Dia89}
\lit{diamond\_1989\_congruence-modules-associated-to-lambda-adic-forms}{\S3: Lemma 3.2, p.~57; Prop.~3.3, p.~58--59}) \stcited:
let $A,B\subseteq V_N(K)$ be the $q$-old subspace (image of $\alpha$) and its orthogonal complement (the $q$-new subspace) for Hida's pairing,
perfect as $p\nmid M(k-2)!$, and $\Omega:=\big((A+V_N(\cO))/V_N(\cO)\big)\cap\big((B+V_N(\cO))/V_N(\cO)\big)\subseteq V_N(K/\cO)$. Then $\Omega$ is a
module for the weight-$k$ Hecke algebra $\mathcal H_Z$ of level $Mq$ (generated by all $T_n$, so containing $U_q$) on which $\fb$ (killing
$A$) and $\fa$ (killing $B$) act by $0$; by (B1), $\Omega\cong\ker(\beta\circ\alpha)$ $\mathcal H_Z$-linearly, where $\mathcal H_Z$ acts on
$W(K/\cO)^2$ by the level-$M$ operators on the diagonal and $U_q$ by the companion matrix of $X^2-T_qX+q^{k-1}$; and
$U_q^2-S_q=(\text{automorphism})\circ\beta\circ\alpha$, so $\Omega\cong\ker(U_q^2-S_q\mid W(K/\cO)^2)$, with $S_q=q^{k-2}$ on the trivial character.
The trivial-character part is a direct summand because $p\nmid\varphi(N)$ (H0) \cite[p.~208]{Dia91}. ($\beta\circ\alpha$ is Diamond's form of
$\Psi^*\Psi$.)
(B3) \emph{Non-Eisenstein localisation.} After localisation at the non-Eisenstein $\m$, $H^1$ and parabolic $H^1_P$ agree (the boundary
cohomology is Eisenstein), so $W(\cO)_\m=H^1(\Gamma,V)_\m\cong T^2$. \stknown{} (standard; no source in \pathtt{literature/} located, open point (N5); \sttoprove{} if a reference is required, and \pathtt{prop:s06-betti}(b) is \stpmc{} modulo this input).
(B4) $T_q^2-(q+1)^2q^{k-2}$ is a non-zero-divisor on $T\otimes K$: an eigenform $g$ of level $M$ with $a_q(g)=\pm(q+1)c$ would have a Frobenius
eigenvalue $\pm c$ at $q$ of absolute value $q^{(k-2)/2}\ne q^{(k-1)/2}$, contradicting the Ramanujan--Petersson bound \cite[Thm~5.6]{Del69}
\lit{deligne\_1969\_formes-modulaires-et-representations-l-adiques}{\S5, Thm 5.6 (Weil $\Rightarrow$ Ramanujan) and the remark on
Petersson's generalisation, p.~19--20 (PDF)} \stcited{} (the Weil conjectures themselves are not in \pathtt{literature/}).

\begin{proof}[Proof of (b)]
Write $T=\T_k(M)_\m$, $\T_k:=\T_k(N)_{\m_N}$ and $\pi'':\T_k\to\T''_k$. If $\T''_k=0$ the left side is $0$; then no root of
$X^2-T_qX+q^{k-1}$ is $\equiv\epsilon c$ mod $\m$, i.e.\ $T_q\not\equiv\epsilon c(q+1)$, and the right side is $0$ as well. So assume $\T''_k\ne0$.

\emph{Old quotient.} The roots of $X^2-T_qX+q^{k-1}$ mod $\m$ are $\epsilon c$ and $q^{k-1}/(\epsilon c)=\epsilon qc$, distinct since
$q\not\equiv1$ (H0). As in Lemma~\ref{lem:s06-decomp}(b) (fixed weight; the role of the $q$-new form is played by $U_q\equiv\epsilon c$, which
gives $\bar U_q^2=q^{k-2}$), $\T''_k\cong T$ with $U_q\mapsto\alpha_q$, the root $\equiv\epsilon c$, and $T_q\equiv\epsilon c(q+1)\bmod\m$. Put
$\beta_q:=q^{k-1}\alpha_q^{-1}\equiv\epsilon qc$. Then
\[T_q-\epsilon(q+1)c=(\alpha_q-\epsilon c)(1-\epsilon qc\,\alpha_q^{-1}),\qquad \alpha_q^2-q^{k-2}=(\alpha_q-\epsilon c)(\alpha_q+\epsilon c),\]
and $1-\epsilon qc\alpha_q^{-1}\equiv1-q$, $\alpha_q+\epsilon c\equiv2\epsilon c$ are units. Hence
$(T_q-\epsilon(q+1)c)=(\alpha_q-\epsilon c)=(\alpha_q^2-q^{k-2})$ as ideals of $T$.

\emph{$\supseteq$.} For a $q$-new eigenform $g$ of weight $k$, level dividing $N$, trivial character, $a_q(g)^2=q^{k-2}$ ((I2) of
Lemma~\ref{lem:s06-decomp}). So $U_q^2-q^{k-2}\in\fa$, and $\pi''(U_q^2-q^{k-2})=\alpha_q^2-q^{k-2}$. Thus $\pi''(\fa)\supseteq(\alpha_q^2-q^{k-2})$.

\emph{$\subseteq$.} By (B2), $\Omega_{\m_N}\cong\ker\big(U_q^2-q^{k-2}\mid W(K/\cO)^2_{\m_N}\big)$. The action of $\mathcal H_Z$ on $W^2$ factors through
$T[X]/(X^2-T_qX+q^{k-1})\cong T\times T$ ($X\mapsto(\alpha_q,\beta_q)$, Hensel as in Lemma~\ref{lem:s06-decomp}(b)); localising at $\m_N$
($U_q\equiv\epsilon c$) keeps the first factor, on which $U_q=\alpha_q$. With (B3),
$W(K/\cO)^2_{\m_N}\cong W(\cO)_\m\otimes_\cO K/\cO\cong(T\otimes_\cO K/\cO)^2$ as $T$-modules, $T$ acting on $\Omega_{\m_N}$ through
$\T''_k\cong T$. For $x\in T$ a unit of $T\otimes K$ (a non-zero-divisor), the snake lemma for $0\to T\to T\otimes K\to T\otimes K/\cO\to0$
gives $\ker(x\mid T\otimes K/\cO)\cong T/xT$. By (B4) this applies to $x=\alpha_q^2-q^{k-2}$ (a unit multiple of $T_q-\epsilon(q+1)c$). Hence
$\Omega_{\m_N}\cong(T/(\alpha_q^2-q^{k-2}))^{2}$, whose annihilator in $T$ is $(\alpha_q^2-q^{k-2})$. Since $\fa$ kills $\Omega$,
$\pi''(\fa)\subseteq\Ann_T(\Omega_{\m_N})=(\alpha_q^2-q^{k-2})$.

\emph{Conclusion.} $\pi''(\fa)=(T_q-\epsilon(q+1)c)$, and $\T_k/(\fa+\fb)\cong\T''_k/\pi''(\fa)\cong T/(T_q-\epsilon(q+1)c)$.
\end{proof}
\noindent It is consistent with Lemma~\ref{lem:s06-decomp}(b): there $\T''_k\cong\T_k(M)_\m$ with $U_q\mapsto\alpha_q$, and
$T_q-\epsilon(q+1)q^{(k-2)/2}=(\alpha_q-\epsilon q^{(k-2)/2})(1-\epsilon q^{k/2}\alpha_q^{-1})$, where the second factor is
$\equiv1-q$ mod $\m$, a unit since $q\not\equiv1$. \provisional{F2: consistency of (b) with the corrected old quotient checked and
recorded; before: no remark.} [Ri84] (the classical source of (b) at $k=2$) is in \pathtt{literature/} only as an image scan
(\pathtt{ribet\_1984\_congruence-relations-between-modular-forms\_scan}) and is not used; \cite[Thm~4A--4C, p.~61--63, Cor.~6.9,
p.~80]{Dia89} \lit{diamond\_1989\_congruence-modules-associated-to-lambda-adic-forms}{Thm 4A p.~61, Cor.~6.9 p.~80} gives the existence of
the congruences ($k=2$, resp.\ ordinary). Over $K/\cO$, $\Omega_\m$ and $\coker(\Psi^*\Psi)_\m$ are both $(T/\delta)^2$: the factor $2$ of (a) is
the rank of $H$ over $T$.

\begin{proof}[Proof of (c)]
Let $\T'_k$ be the $q$-new quotient of $\T_k=\T_k(N)_{\m_N}$, $\cc=\mathrm{im}(\fb\to\T'_k)$, and $\lambda'_f:\T'_k\to\cO$ the eigensystem of $f$.
By Def.~\ref{def:s06-ctam} at weight $k$, $\T'_k/\cc\cong\T_k/(\fa+\fb)$, and $\lambda'_f(\cc)=\cc_q(f)$. Since $\lambda'_f$ is surjective and maps $\cc$
onto $\cc_q(f)$, it induces a surjection $\T'_k/\cc\twoheadrightarrow\cO/\cc_q(f)$, so
\[\ord_\varpi\cc_q(f)\le\mathrm{length}_\cO\T_k/(\fa+\fb)=\mathrm{length}_\cO T/(\delta)=\tfrac12\,\mathrm{length}_\cO\coker(\Psi^*\Psi)_\m\]
by (b) and (a). If $\rk_\cO\T'_k=1$, then $\T'_k$ is an $\cO$-torsion-free finite $\cO$-algebra with $\T'_k\otimes K=K$, so $\T'_k=\cO$ (an
$\cO$-order in $K$ containing $1$) and $\lambda'_f$ is an isomorphism; the surjection above is then an isomorphism and equality holds.
\end{proof}

\begin{remark}[weights $k\ge p$]\label{rem:s06-betti-kp}
For $q\equiv-1$, (a) fails: the cokernel is no longer $(\T/\delta)^2$, and the two sign factors of $\delta$ lie in $\m$ both or neither (they differ by $2(q+1)c\in\m$). For $k\ge p$
(T4 beyond $k=12$) \pathtt{if:s03-amalgam} is outside its range $l>k$, and (B1) needs $p\nmid(k-2)!$. One needs the amalgam version
\emph{localised at the non-Eisenstein $\m$ and after the ordinary projector}. There the end terms $H^0$ are Eisenstein and vanish
(Risk R3). Stage 03 does not export that version, so every use of Prop.~\ref{prop:s06-betti} for $k\ge p$ (T4 at $k=28,44,60,124$)
is \stexp. \provisional{F19: uses of the amalgam for $k\ge p$ marked \stexp; before: no status for the $k\ge p$ uses.}
\end{remark}
